%!TEX root = main.tex
\section{Discussion}
\label{sec:discussion}

\subsection{What the Hybrid Is Good For}
\label{sec:value}

The result is better read as a successful audit: 38 of 6{,}680 decisions change,
each described by a channel, a direction and a mechanism.  Moreover, the
end-to-end bilingual evaluation demonstrates deployment viability (87.10\%
accuracy, 9.07 ms latency, 110.2 texts/s throughput), showing that a hybrid architecture can
be both explainable and production-ready.

\emph{Triage and accountability.}  A queue can sort comments by whether the
decision was rule-supported: a \hate{} decision with a matched span and
$r_H=0.8$ is easier to review than one with no lexical evidence, and the span,
the entry and the arithmetic are what a platform can defend in policy terms.

\emph{Editing and diagnosis.}  Adding a slur variant is a configuration change,
not a retraining run, and an edit affects the \hate{} channel (weight 0.8) far
more than the others.  The decomposition exists only because the channels are
separate: a monolithic classifier improving by 0.0027 macro-F1 would give no
handle on whether the change was lexical or neural.

\subsection{Why Accuracy Does Not Move (Much)}
\label{sec:why_neutral}

First, \emph{coverage is low}: the rule fires on 15.75\% of comments, so a
perfect rule channel could still only change decisions in that subset.  Second,
\emph{the majority class dominates}: with 83\% \clean{}, only the 1{,}132 toxic
comments can improve, and 38 changes is 3.4\% of that support---enough to matter
operationally, too little to move a macro-average beyond its confidence
interval.  Third, \emph{tuning is safety-oriented}: protecting \clean{} and
gating \hate{} on lexical support produces both the de-flagging and the missed
corrections.  Neutrality is a \emph{choice}, not an absence of signal.

\subsection{The Bilingual Asymmetry}
\label{sec:asymmetry_bilingual}

The mixed evaluation reveals a clear language gap: HATE F1 VI 0.6051 vs EN
0.4537 (Vietnamese lexicon has 594 HATE entries versus English 86; Davidson
training only 5.8\% HATE), OFFENSIVE F1 EN 0.9270 vs VI 0.4142 (English
backbone aligns with Davidson's 77.8\% OFFENSIVE majority; ViHSD only 6.65\%
OFFENSIVE), CLEAN F1 VI 0.9344 vs EN 0.8541 (English experiences mild
over-flagging on clean text). This asymmetry is data-driven: lexicon extraction
from corpora with different class priors produces different coverage.

\subsection{Design Implications}
\label{sec:implications_bilingual}

Five changes follow:
\begin{enumerate}
    \item \textbf{Evidence-gated fusion}: when the rule fires strongly against
          a high-confidence \clean{} prediction, abstain or escalate.
    \item \textbf{Separate thresholds per mechanism}: requiring $r_H\ge0.5$
          before the \hate{} channel can override would make the rule-driven
          and threshold-driven paths separately tunable.
    \item \textbf{Per-class calibration} of the \clean{} channel, which currently
          blocks the missed corrections.
    \item \textbf{Balanced English lexicon}: expand HATE entries by merging
          HateXplain/Gab sources and retrain the English rule engine.
    \item \textbf{Multilingual fallback checkpoint}: for \texttt{mix}/\texttt{multilingual}
          samples (5.6\%), route to a multilingual mDeBERTa instead of rule-only.
\end{enumerate}

\subsection{Practical Deployment}
\label{sec:deployment}

The latency profile (9.07 ms mean, 110.2 texts/s throughput) supports synchronous
moderation for text comments.  Routing (5.22 ms) and AI inference (3.85 ms) are the
main components; P95 latency is 12.69 ms.  Tail latency can reach up to 2.1 s during cold-start tokenizer
initialization; this is a one-time cost that can be pre-warmed.  A production
deployment should pre-load both checkpoints and the two rule engines at
startup, cache tokenizer initialization, and set $\alpha$ per language from
validation-frozen weights.